\chapter{Multi-Socket Machines and Interconnects}\label{ch:ll:multi-socket-machines-and-interconnects}

A trading server with two processor sockets, a network card in a slot wired to the second socket, and a strategy
thread that the operating system happened to start on the first: every packet the card delivered was written near
socket 1 and read from socket 0, and every order went back the other way. Each crossing costs roughly a hundred
nanoseconds more than staying on one socket, several times per message, on a path whose whole software budget is a
few microseconds. Nobody had chosen the placement; nobody had looked at the machine's topology. This chapter describes
the machine above the core: memory attached to sockets, the links between sockets, the bus that connects the network
card, and how to read all of it from the operating system and place a trading process on it.

\section{Non-uniform memory access}

\begin{definition}[Non-uniform memory access, NUMA node]\label{def:ll:multi-socket-machines-and-interconnects:numa}
In a machine with \emph{non-uniform memory access}\index{non-uniform memory access} (NUMA) each processor socket has its
own memory controllers and memory; a core reaches its own socket's memory faster than another socket's, which it
reaches over the \omterm{def:ll:multi-socket-machines-and-interconnects:link}{socket interconnect}. A \emph{NUMA node}\index{NUMA node} is a set of CPUs and the memory close to them,
usually one socket (or part of one).
\end{definition}

\begin{definition}[First-touch allocation]\label{def:ll:multi-socket-machines-and-interconnects:first}
Under \emph{first-touch allocation}\index{first-touch allocation}, the default policy of Linux, a page of memory is
placed on the \omterm{def:ll:multi-socket-machines-and-interconnects:numa}{NUMA node} of the CPU that first writes it (the page fault that creates it), not the one that allocated
the address range.
\end{definition}

First touch makes placement a question of who initialises what. A trading process that builds its order books, rings
and pools in its main thread during start-up, before creating its threads, puts all of them on the main thread's node,
wherever the hot threads later run. The rule: pin each hot thread first (chapter~9), then let it allocate and touch its
own data; or bind the memory explicitly (\texttt{numactl --membind}, \texttt{mbind}).

\begin{example}[A book allocated on the wrong node]\label{ex:ll:multi-socket-machines-and-interconnects:touch}
The main thread, on node 0, allocates and zeroes 64~MiB of order books with \omterm{def:ll:memory-hierarchy-and-caches:tlb}{huge pages} at start-up, then starts the
feed thread on node 1. Every book update now misses to the other socket's memory. The fix is one line: move the
initialisation into the feed thread, after it has pinned itself, or pass the node to the allocator.
\end{example}

\section{The socket interconnect and remote caches}

\begin{definition}[Socket interconnect]\label{def:ll:multi-socket-machines-and-interconnects:link}
The \emph{socket interconnect}\index{socket interconnect} is the set of point-to-point links between the processor
sockets of a machine, which carry remote memory accesses and the cache-coherence traffic between sockets.
\end{definition}

Every coherence transaction of chapter~3 that involves a line held on the other socket crosses the interconnect. A
published measurement of a two-socket server of a recent Intel generation found that bouncing a \omterm{def:ll:memory-hierarchy-and-caches:line}{cache line} between two
cores took 59 nanoseconds on average within a socket (81 at worst) and 138 nanoseconds across sockets; older
generations were similar, up to 150 across. A hand-off between two threads, the feed handler publishing a message and
the strategy reading it, is exactly such a bounce (chapter~12): placed on two sockets, it more than doubles.

The same measurement on this book's laptop, which has one socket and runs under a hypervisor, is
\cref{fig:ll:multi-socket-machines-and-interconnects:c2c}. It has structure: pairs of neighbouring virtual CPUs hand a
line over in about 40 nanoseconds, most other pairs in 170 to 200. The laptop's physical cores sit in clusters sharing
a second-level cache, and the virtual CPUs are mapped onto them by the hypervisor, not visibly to the guest; the
matrix, not the CPU numbering, is what tells which threads should talk to each other.

\begin{omfigure}
\begin{tikzpicture}
\begin{axis}[width=8.6cm,height=8.6cm,view={0}{90},xmin=-0.5,xmax=21.5,ymin=-0.5,ymax=21.5,y dir=reverse,
  xtick={0,3,6,9,12,15,18,21},ytick={0,3,6,9,12,15,18,21},xlabel={CPU},ylabel={CPU},
  tick label style={font=\small},label style={font=\small},
  colormap={om}{color(0)=(white) color(0.3)=(omDef!30) color(1)=(omDef)},
  colorbar,colorbar style={ylabel={ns, one way},ylabel style={font=\small},tick label style={font=\footnotesize}},
  point meta min=0,point meta max=200,table/col sep=comma]
\addplot[matrix plot*,mesh/cols=22,point meta=explicit] table[x=b,y=a,meta=ns]{figdata/low-latency/04-multi-socket-machines-and-interconnects/measured_c2c.csv};
\end{axis}
\end{tikzpicture}

\omcaption{One-way latency of handing a \omterm{def:ll:memory-hierarchy-and-caches:line}{cache line} between two threads pinned to two of the laptop's 22 virtual CPUs
(20\,000 round trips, best of three). Neighbouring CPUs are fast; the rest are four to five times slower. Measured on a
laptop (Intel Core Ultra 7 155H) under WSL2, no isolated cores. Data: \texttt{bench\_c2c.py}.}\label{fig:ll:multi-socket-machines-and-interconnects:c2c}
\end{omfigure}

\section{The peripheral bus and direct memory access}

\begin{definition}[PCI Express, direct memory access]\label{def:ll:multi-socket-machines-and-interconnects:pcie}
\emph{PCI Express}\index{PCI Express} (PCIe) is the serial point-to-point bus that connects devices (network cards,
storage, accelerators) to a processor socket, in lanes grouped by 1 to 16. In \emph{direct memory access}\index{direct
memory access} (DMA) a device reads and writes host memory itself, without the processor copying each byte: a network
card writes each received packet into a buffer the driver gave it and signals its arrival by a descriptor, an interrupt
or a flag the program polls.
\end{definition}

\begin{definition}[Direct cache access]\label{def:ll:multi-socket-machines-and-interconnects:dca}
\emph{Direct cache access}\index{direct cache access} lets a device's DMA writes land in the processor's last-level cache
instead of memory, so that the core that processes the data finds it in cache; Intel's implementation (Data Direct I/O)
is transparent to software.
\end{definition}

Every PCIe slot belongs to one socket: its lanes come from that socket's processor. A card's packets therefore arrive in
that socket's cache or memory, and the card's registers, which the program writes to send, are across the interconnect
from any core on the other socket. Each transaction on the bus adds its own latency, which is why low-latency network cards and their user-space drivers
minimise the number of bus transactions per packet (One Quant Book~14, chapter~3).

\begin{dated}{2026-09}{PCI Express generations}\label{dat:ll:multi-socket-machines-and-interconnects:pcie}
PCIe 4.0 signals at 16~GT/s per lane and PCIe 5.0, whose specification PCI-SIG released in May 2019, at 32~GT/s; both
encode 128 bits in 130. A 16-lane PCIe 5.0 link therefore carries $32 \times 16 \times 128/130/8 \approx 63$~GB/s in each
direction, far more than a 100~Gb/s network port needs (12.5~GB/s): for a trading card the bus matters for its latency,
not its bandwidth.
\end{dated}

\section{Where the network card sits}

\begin{definition}[Device locality]\label{def:ll:multi-socket-machines-and-interconnects:locality}
The \emph{device locality}\index{device locality} of a thread or a buffer is its \omterm{def:ll:multi-socket-machines-and-interconnects:numa}{NUMA node} relative to the device it
exchanges data with: local when they are on the same node, remote otherwise.
\end{definition}

\begin{omfigure}
\begin{tikzpicture}[font=\footnotesize,
  sock/.style={draw=omDef,fill=omDef!6,rounded corners=3pt,minimum width=4.2cm,minimum height=2.4cm},
  mem/.style={draw=gray,fill=gray!12,minimum width=4.2cm,minimum height=0.55cm},
  thr/.style={draw=omThm,fill=omThm!12,rounded corners=2pt,minimum width=1.6cm,minimum height=0.5cm}]
\node[mem] (m0) at (0,2.1) {memory of node 0};
\node[mem] (m1) at (7,2.1) {memory of node 1};
\node[sock] (s0) at (0,0) {};
\node[sock] (s1) at (7,0) {};
\node[anchor=north] at (s0.north) {socket 0: cores, caches};
\node[anchor=north] at (s1.north) {socket 1: cores, caches};
\node[thr] (bad) at (0,-0.35) {strategy (remote)};
\node[thr,fill=omMeth!15,draw=omMeth] (good) at (7,-0.35) {strategy (local)};
\draw[thick] (m0) -- (s0); \draw[thick] (m1) -- (s1);
\draw[<->,very thick,omProp] (s0.east) -- node[above,align=center]{interconnect\\ 138 ns a line} (s1.west);
\node[draw=black,fill=gray!20,minimum width=2.4cm,minimum height=0.6cm] (nic) at (7,-2.3) {network card};
\draw[<->,thick] (nic) -- node[right]{PCIe lanes of socket 1} (s1.south);
\end{tikzpicture}

\omcaption{A two-socket server whose network card is wired to socket 1. A strategy thread on socket 1 reads packets from
its own cache; one on socket 0 fetches every packet line, and sends every order, across the interconnect, at about 138
nanoseconds a line against 59 within a socket (published averages for a recent generation).}\label{fig:ll:multi-socket-machines-and-interconnects:server}
\end{omfigure}

\begin{proposition}[Cost of remote placement]\label{prop:ll:multi-socket-machines-and-interconnects:remote}
If a message brings $k_{\mathrm{in}}$ \omterm{def:ll:memory-hierarchy-and-caches:line}{cache lines} from the card to the hot thread and an order sends $k_{\mathrm{out}}$ back,
each transferred one after the other, and a line costs $c_{\mathrm L}$ within the socket and $c_{\mathrm R}$ across it, then
remote placement adds $(k_{\mathrm{in}} + k_{\mathrm{out}})(c_{\mathrm R} - c_{\mathrm L})$ to the \omterm{def:ll:where-latency-comes-from:t2t}{tick-to-trade latency}.
\end{proposition}

\begin{proof}
Each of the $k_{\mathrm{in}} + k_{\mathrm{out}}$ dependent transfers costs $c_{\mathrm R}$ instead of $c_{\mathrm L}$; the rest
of the path is unchanged.
\end{proof}

With two lines each way and the published 59 and 138 nanoseconds, the proposition gives 316 nanoseconds per message:
more than a tenth of the \qty{2.5}{\micro\second} a very good software system takes wire to wire (chapter~1). The
placement rules follow: the card, the hot threads and their memory on one node; each hot thread on its own physical core,
its hyper-threading sibling idle; the interrupts of other devices and the operating system's work on the other node or
on housekeeping cores (chapter~13); and every thread that exchanges data with the hot path per message (the risk gate,
the logger's ring) on the same node.

\section{Reading a machine's topology}

Linux exposes the topology in sysfs. Each CPU's physical package, core and hyper-threading siblings are under
\path{/sys/devices/system/cpu/cpu*/topology/}; each node's CPUs under \path{/sys/devices/system/node/node*/cpulist};
each PCI device's node in \path{/sys/bus/pci/devices/*/numa_node}, which holds $-1$ when the firmware did not say; and
the device behind a network interface is the link \path{/sys/class/net/<name>/device}. Tools such as
\texttt{lscpu}, \texttt{numactl --hardware} and \texttt{lstopo} print the same facts; a placement plan should read them
itself, so that a moved card or a new server cannot silently invalidate it.

This laptop shows what can go wrong. Linux sees one node, 22 CPUs described as eleven cores of two threads (the processor
has 6 two-thread performance cores and 10 single-thread efficient cores), and a network interface that is a synthetic
adapter of the hypervisor, not a PCI device: there is no card to be local to. The build of this chapter therefore
treats an unknown node as a fact to report, not to guess.

\section{Tutorial: a placement plan and the core-to-core matrix}\label{tut:ll:multi-socket-machines-and-interconnects:tut}

\begin{tutorial}
\textbf{Goal.} Read a machine's topology, place the hot threads of a trading process on it, check the plan, and measure
the cost of a hand-off between every pair of CPUs. \textbf{End state:} a checked plan for a two-socket fixture and for
this laptop, and \cref{fig:ll:multi-socket-machines-and-interconnects:c2c}.
\begin{tutsteps}
\item \textbf{Read the tree.} \texttt{read\_topology(root)} parses the sysfs files above; the tests run it on two
synthetic trees written by \texttt{make\_coreplan\_fixture.py} (a two-socket, 32-CPU server with its card on node 1, and a
one-node machine whose card's node is unknown).
\item \textbf{Place.} The hot roles take the highest-numbered physical cores of the card's node, one each, leaving their
siblings idle; CPU 0's core and every other CPU go to housekeeping.
\omcode{code/firm/coreplan/firm_coreplan.py}{115}{131}{A placement plan from the topology.}\label{lst:ll:multi-socket-machines-and-interconnects:plan}
\item \textbf{Check.} \texttt{check(plan, topo)} lists every violated rule: a hot thread on the wrong node, two hot
threads on one physical core, a busy sibling, CPU 0 not in housekeeping.
\item \textbf{Measure hand-offs.} Two pinned threads bounce one padded atomic flag; the time per one-way transfer is the
core-to-core latency. \texttt{python bench\_c2c.py} runs it for every pair.
\omcode{code/low-latency/04-multi-socket-machines-and-interconnects/cpp/ll_c2c.hpp}{23}{50}{Bouncing a cache line between two pinned threads.}\label{lst:ll:multi-socket-machines-and-interconnects:pingpong}
\end{tutsteps}
\textbf{What to change next.} Add a rule to \texttt{check} that the logger is on the card's node too, and decide whether
it should be; rerun the ping-pong with the flag and a payload of four lines, and compare the cost per line.
\end{tutorial}

\section{Build: the core placement plan}\label{bld:ll:multi-socket-machines-and-interconnects:coreplan}

\begin{build}
\bfield{Purpose} One description of where every thread of the trading process runs, derived from the machine rather than
typed by hand: chapter~9 pins threads from it, chapter~13 audits a host against it, chapter~26 records it with every
measurement.
\bfield{Interface} Python \texttt{firm\_coreplan}: \texttt{read\_topology(root)} returning \texttt{Topology} (CPUs with
package, core, node and siblings; nodes; PCI cards with node and local CPUs; interfaces), \texttt{Topology.card\_for(ifname)},
\texttt{plan(topo, ifname, hot, cold)} returning \texttt{Plan(node, roles, idle, housekeeping)}, \texttt{check(plan,
topo)}, \texttt{remote\_penalty\_ns}.
\bfield{Rules} Hot roles on the card's node, one physical core each, siblings idle, never CPU 0's core; an unknown node
($-1$) falls back to the lowest node and is reported; the plan fails loudly when the node has too few cores.
\bfield{Acceptance tests} \texttt{code/firm/coreplan/tests/}: parsing CPU lists; the two-socket fixture's topology and a
sound plan on node 1; plans with a thread on the wrong socket and two threads on one core caught; the one-node fixture with
an unknown card.
\bfield{Stretch} Read cache topology (\path{/sys/devices/system/cpu/cpu*/cache/}) and place threads that exchange data per
message on cores sharing a cache level; include the measured core-to-core matrix in the plan.
\end{build}

\begin{omsources}
\item Chips and Cheese, ``Core to core latency data on large systems'', 2023.
\item Linux kernel documentation, ``NUMA memory policy''; ABI documentation, \texttt{sysfs-bus-pci}.
\item Intel, ``Intel Data Direct I/O technology''.
\item PCI-SIG, PCI Express 5.0 specification announcement, 2019.
\end{omsources}

\section{Exercises}

\begin{exercise}[$\star$]\label{exo:ll:multi-socket-machines-and-interconnects:1}
What is the usable bandwidth in each direction of an 8-lane PCIe 4.0 link? Is it enough for two 25~Gb/s ports?
\end{exercise}

\begin{exercise}[$\star$]\label{exo:ll:multi-socket-machines-and-interconnects:2}
A packet of 300 bytes lands in the cache of the card's socket. How many \omterm{def:ll:memory-hierarchy-and-caches:line}{cache lines} must a core on the other socket fetch to
read it?
\end{exercise}

\begin{exercise}[$\star$]\label{exo:ll:multi-socket-machines-and-interconnects:3}
\path{/sys/bus/pci/devices/0000:81:00.0/numa_node} contains 1 and the strategy runs on CPU 5, whose node is 0. What
does \cref{prop:ll:multi-socket-machines-and-interconnects:remote} give with three lines in, one out, and the published
figures?
\end{exercise}

\begin{exercise}[$\star\star$]\label{exo:ll:multi-socket-machines-and-interconnects:4}
The feed handler, pinned to node 1, publishes each message into a ring that the main thread allocated and zeroed at
start-up on node 0. Where does each write go, and how do you fix it?
\end{exercise}

\begin{exercise}[$\star\star$]\label{exo:ll:multi-socket-machines-and-interconnects:5}
Read \cref{fig:ll:multi-socket-machines-and-interconnects:c2c}: which pairs of CPUs would you give to a feed thread and a
strategy thread that hand over every message, and what would a hand-off cost between the pair you chose and a random pair?
\end{exercise}

\begin{exercise}[$\star\star$]\label{exo:ll:multi-socket-machines-and-interconnects:6}
Why must a hot thread's hyper-threading sibling stay idle, and what does the plan lose by it on the two-socket fixture?
\end{exercise}

\begin{exercise}[$\star\star\star$]\label{exo:ll:multi-socket-machines-and-interconnects:7}
\emph{Coding.} Extend \texttt{firm\_coreplan.plan} to place four hot threads on a node with only three usable physical cores
by sharing one core between the two threads that exchange the least data. Which pair would you choose in a feed, book,
strategy and gateway pipeline, and what does \texttt{check} need to accept?
\end{exercise}

\begin{exercise}[$\star\star\star$]\label{exo:ll:multi-socket-machines-and-interconnects:8}
\emph{Find the flaw.} ``We pinned our strategy with \texttt{taskset -c 5} in the service's start script, so it is on the
right socket.''
\end{exercise}

\section{Problem: Wrong Socket}

\begin{problem}[{Weekend problem --- a server nobody had looked at}]\label{pb:ll:multi-socket-machines-and-interconnects:1}
A firm's new two-socket server has its network card on node 1. The feed handler, strategy and gateway were started by a
script without any pinning, and the scheduler put them on node 0. Each market-data message brings two \omterm{def:ll:memory-hierarchy-and-caches:line}{cache lines} to the
strategy; each order sends two lines to the card. Take the published figures: 59 nanoseconds per line within a socket, 138
across.

\medskip\noindent\textbf{Part I --- The topology.}
\begin{enumerate}
\item Which sysfs files tell the node of the card and of each CPU?
\item On the fixture \texttt{two\_socket}, which CPUs are on node 1, and what are the siblings of CPU 17?
\item What does \texttt{plan} return for the three hot roles, and which CPUs does it leave idle?
\item Why does the plan never use CPU 0's core?
\end{enumerate}

\medskip\noindent\textbf{Part II --- The penalty.}
\begin{enumerate}[resume]
\item Apply \cref{prop:ll:multi-socket-machines-and-interconnects:remote}: what does the wrong socket add per message?
\item What fraction of a \qty{2.5}{\micro\second} wire-to-wire budget is that?
\item If the strategy also reads its order book, allocated by the main thread on node 0, what happens when the threads are
moved to node 1 without moving the book?
\item Which of the two mistakes is worse, and why?
\end{enumerate}

\medskip\noindent\textbf{Part III --- Hand-offs.}
\begin{enumerate}[resume]
\item The feed thread hands every message to the strategy through a ring: one line each way. What does the hand-off cost
within a socket and across?
\item On this laptop, what does it cost between neighbouring CPUs and between others (\cref{fig:ll:multi-socket-machines-and-interconnects:c2c})?
\item How would you choose the two CPUs from the matrix alone?
\item Why can the matrix change from one boot of a virtual machine to the next?
\end{enumerate}

\medskip\noindent\textbf{Part IV --- The fix.}
\begin{enumerate}[resume]
\item State the \emph{named result}: the latency added per message by the wrong socket, in nanoseconds and as a share of
a \qty{2.5}{\micro\second} budget.
\item Write the steps of the fix, in order.
\item How do you verify after the fix that each thread runs where the plan says?
\item What should happen when a hardware change moves the card to node 0?
\item Which other devices' interrupts should be moved away from the hot cores?
\item What does \texttt{numa\_node} $= -1$ mean, and what should the plan do?
\item What would you record with every latency measurement so that a placement change is visible later?
\item In one sentence: why is placement a latency decision?
\end{enumerate}
\end{problem}

\section{Interview questions}

\begin{interviewq}[$\star$ \iqroles{developer}]\label{iq:ll:multi-socket-machines-and-interconnects:1}
What is NUMA, and what does it mean for a latency-sensitive process?
\end{interviewq}

\begin{interviewq}[$\star\star$ \iqroles{developer}]\label{iq:ll:multi-socket-machines-and-interconnects:2}
What is \omterm{def:ll:multi-socket-machines-and-interconnects:first}{first-touch allocation}, and how can it put a thread's memory on the wrong node?
\end{interviewq}

\begin{interviewq}[$\star\star$ \iqroles{developer}]\label{iq:ll:multi-socket-machines-and-interconnects:3}
How do you find which socket a network card is attached to on Linux?
\end{interviewq}

\begin{interviewq}[$\star\star$ \iqroles{developer}]\label{iq:ll:multi-socket-machines-and-interconnects:4}
Two threads exchange a message a microsecond. Where would you place them on a two-socket machine, and why?
\end{interviewq}

\begin{interviewq}[$\star\star$ \iqroles{developer}]\label{iq:ll:multi-socket-machines-and-interconnects:5}
What is DMA, and what does \omterm{def:ll:multi-socket-machines-and-interconnects:dca}{direct cache access} change for a packet's path to the application?
\end{interviewq}

\begin{interviewq}[$\star\star\star$ \iqroles{developer}]\label{iq:ll:multi-socket-machines-and-interconnects:6}
Design a start-up procedure for a trading process that guarantees correct placement of its threads and memory on any
server in the fleet.
\end{interviewq}
